\section{Limitations}
\label{sec:limitations}

\paragraph{Target leakage is mitigated, not eliminated}
The 8/8 regime reconstructs the index formula and we report it as a benchmark rather
than a prediction result. Even at 7/8 and 6/8, the retained measurements are formula
contributors, so part of the remaining performance is partial formula reconstruction
rather than environmental inference. A fully independent regime -- predicting the index
from metadata and spatiotemporal context alone, with no formula contributors -- would
bound this, and we did not run it.

\paragraph{A point mass at the ceiling is reported as one number}
With 35.92\% of records at the upper bound, pooled $R^2$ and RMSE conflate two tasks:
identifying records that are fully guideline-compliant, and estimating how far the
remainder fall short. The figures reported here are pooled. Metrics separated by ceiling
status would characterise the models more honestly, particularly in the low-availability
regimes where identifying compliance is likely the easier half of the problem. We note
that this is a bounded target rather than a censored one
(Section~\ref{sec:ceiling}): no latent value exists above 100, so the remedy is
stratified reporting, not a censored-regression estimator.

\paragraph{Imputation cannot be undone}
Up to 31.6\% of individual measurement columns were median-filled before publication
(Section~\ref{sec:imputation}). The original missingness pattern is unrecoverable, so the
mechanism cannot be classified as MCAR, MAR or MNAR. Because the fill is concentrated in
England, fill status, country and target are mutually confounded, and no split design
removes this.

\paragraph{A probable defect in the upstream index}
Dissolved oxygen correlates with the target at $r = -0.031$. Dissolved oxygen is a
``more-is-better'' parameter and would normally be among the most informative inputs. The
upstream implementation appears to treat its guideline as an upper bound to remain below,
inverting its contribution. We report the index as published and flag this rather than
silently correcting it.

\paragraph{Composition and coverage}
Effluent and sewage comprise 22.1\% of records, so the dataset is not purely natural
surface water. England supplies 75.3\% of records and its series begins in 2000, so the
corpus is not geographically or temporally balanced, and aggregate metrics are dominated
by one country.

\paragraph{What the split does and does not establish}
Because whole monitoring stations are allocated to a single partition, every test record
comes from a site the model has never seen, and the reported figures are therefore
estimates of generalisation to new stations. They are not estimates of generalisation
across time: the split is not chronological, so a test site's records may fall in the
same period as the training data, and nothing here speaks to extrapolation into future
years or to conditions outside the observed range. The Chinese subset is a further
exception, being single-site and allocated entirely to training
(Section~\ref{sec:split}); its per-country figures are in-sample.

\paragraph{Unequal training budgets}
Random Forest is fitted on a 100{,}000-record subsample while XGBoost uses the full
1{,}986{,}397-record training partition (Section~\ref{sec:models}). Validation and test
records, preprocessing and feature sets are identical, but the training data is not. The
margin reported in favour of XGBoost therefore confounds algorithmic capability with
training-set size, and should be read as an upper bound on the true algorithmic gap.
Equalising the budgets is the first experiment we would add.

\paragraph{China cannot be evaluated out-of-sample}
China contributes a single monitoring site, which cannot be divided across partitions
without breaking the grouping guarantee and is therefore allocated wholly to training. It
contributes no test records and is absent from Table~\ref{tab:country}. Nothing in this
paper evidences performance on Chinese waters.

\paragraph{Subset selection}
The four availability regimes are interpretable rather than exhaustive. There are eight
ways to withhold one measurement and seventy ways to withhold four; we report one
informative choice at each level rather than averaging over all of them. The degradation
curve should therefore be read as evidence that performance collapses once formula
reconstruction is impossible, not as a precise estimate of expected loss under arbitrary
missingness.
